\documentclass{article}
\usepackage[utf8]{inputenc}
\usepackage[T1]{fontenc}
\usepackage[a4paper]{geometry}		% Reduire les marges
\geometry{hmargin=3cm,vmargin=3.5cm}
\usepackage[english]{babel} 
\usepackage{amsfonts}
\usepackage{amsthm}
\usepackage{amsmath}
\usepackage{amssymb}
\usepackage{hyperref}
\usepackage{array}
\usepackage[svgnames]{xcolor}
\usepackage{xpatch}
\usepackage{tikz}
\usepackage{mathrsfs}
%\usepackage{graphicx}
%\usepackage{caption}
\usepackage[svgnames]{xcolor}
\usepackage{eurosym}
\usetikzlibrary{positioning,calc}
\renewcommand{\thesection}{\Roman{section}}
\xpatchcmd{\proof}{\itshape}{\normalfont\proofnameformat}{}{}
\newcommand{\proofnameformat}{\bfseries}
\theoremstyle{definition}
\newtheorem{defi}{Definition}[section]
\newtheorem{theo}{Theorem}[section]
\newtheorem{prop}[theo]{Proposition}
\newtheorem{lem}[theo]{Lemma}
\newtheorem{coro}[theo]{Corollary}
\newtheorem{exe}{Example}[section]
\newtheorem{rem}{Remark}[section]
\newtheorem{nota}{Notation}[section]
\renewcommand\qedsymbol{$\blacksquare$}
\numberwithin{equation}{section}

\DeclareMathOperator{\R}{\mathbb{R}}
\DeclareMathOperator{\Q}{\mathbb{Q}}
\DeclareMathOperator{\N}{\mathbb{N}}
\DeclareMathOperator{\E}{\mathbb{E}}
\DeclareMathOperator{\C}{\mathcal{C}}

\title{Exo 3, correction}       
\author{Guillaume Woessner}
\date{}                       % sans date : celle du jour de compilation

%\pagestyle{headings}          % Pour mettre des entetes avec les titres
                              % des sections en haut de page
\sloppy                       % Ne pas faire deborder les lignes dans la marge

\begin{document}

\maketitle                    % Faire un titre utilisant les donnees
                              % passees : \title, \author et \date

\begin{abstract}
\begin{center}
Les temps d'arrêt.
\end{center}
\end{abstract}

%\tableofcontents              % Table des matieres


Dans tout ce document, l'espace de probabilité considéré est $(\Omega,\mathcal{F},(\mathcal{F}_n)_n,\mathbb{P})$.



\paragraph{Exercice 1} 
Enoncez le théorème de décomposition de Doob pour les sous-martingales et pour les sur-martingales.\\
Soit $(X_n)_n$ un processus stochastique adapté à la filtration, tel que $\E[|X_n|]<\infty$. Quelle pourrait être une décomposition de Doob pour ce processus ? Montrez qu'effectivement un telle décomposition existe.

\begin{proof}
La même preuve que celle du cours fonctionne pour montrer que $X_n = M_n + A_n$ avec $A_n$ un processus prédictible (\textit{ie} un processus tel que $A_n$ est mesurable par rapport à $\mathcal{F}_{n-1}$) donné par 
$$A_n := \sum_{k=1}^n \E[X_k~|~\mathcal{F}_{k-1}]-X_{k-1},$$
et $M_n$ est une martingale donnée par 
$$M_n = X_n - A_n = X_0 + \sum_{k=1}^n X_k - \E[X_k~|~\mathcal{F}_{k-1}].$$
\end{proof}



\paragraph{Exercice 3} 
Soit $S_n$ une marche aléatoire simple sur $\mathbb{Z}$, et $\tau := \inf_{n\in\N} \{S_n=1\}$. Montrez que $\tau$ est un temps d'arrêt et que $\tau<\infty$ \textit{as}.

\begin{proof}
C'est un temps d'arrêt car 
\begin{align*}
    \{\tau \leq n\} &= \{\exists k\leq n ~:~S_k = 1\}\\
    &= \cup_{k\leq n}  \{ \#\{i\leq k~:~X_i =1\} - \#\{i\leq k ~:~X_i =-1\} = 1\}.
\end{align*}
Or ces évenements sont en nombres finis, et déterminés par les $(X_i ~:~ i\leq n)$. Donc $\{\tau\leq n\}$ est bien mesurable par rapport à $\mathcal{F}_n$.\\
Pour la seconde propriété, voir le corrigé de la Feuille 5 Exercice 1.
\end{proof}



\paragraph{Exercice 5}
Soient $\tau_1$ et $\tau_2$ deux temps d'arrêt. Est-ce que $\tau_1+\tau_2$ est encore un temps d'arrêt ? Pourquoi ?

\begin{proof}
On a 
\begin{align*}
    \{\tau_1+\tau_2 = n\} = \bigcup_{k\leq n}\{\tau_1=k\}\cap\{\tau_2=n-k\}
\end{align*}
Or ces évenements sont en nombres finis, et mesurables par rapport à $\mathcal{F}_n$, ainsi $\{\tau_1+\tau_2 = n\}$ est mesurable par rapport à $\mathcal{F}_n$. Enfin, il suffit de noter que 
$$\{\tau_1+\tau_2 \leq n\} = \bigcup_{k\leq n} \{\tau_1+\tau_2 = k\}.$$
\end{proof}


%\paragraph{Exercice 6}
%Let $\tau_1$ and $\tau_2$ be two stopping times such that $\tau_1 \leq \tau_2$ \textit{as}. Show that $\mathcal{F}_{\tau_1}\subset \mathcal{F}_{\tau_2}$

%\paragraph{Exercice 7}
%Let $(X_n)_n$ be a sequence of \textit{i.i.d.} random variables adapted to the filtration, such that $\E[|X_1|]<\infty$, and $N$ an integrable stopping time.\\
%Give the definition of $Y := \sum_{i=1}^N X_i$, and show that $\E[Y]=\E[N]\E[X_1]$.




\end{document}
